\section{Introduction}

Expansion-based formulations of the Fermi paradox connect technological capability to persistent settlement or self-replicating activity. \citet{Hart1975} and \citet{Tipler1980} gave influential versions of this argument; \citet{Brin1983} reviewed the broader assumptions behind the resulting ``Great Silence.'' Diffusion, finite settlement lifetimes, stellar motion, and heterogeneous sites complicate the route from launch capability to complete occupation \citep{NewmanSagan1981,CarrollNellenback2019}. Cosmological models nevertheless show why even a small number of persistently expanding lineages can matter \citep{Olson2015,Hanson2021,Kipping2026}. This paper examines a behavioural premise rather than denying these propagation arguments: does a civilization that can expand necessarily continue to find additional external investment worthwhile?

The computational inward transition hypothesis (\CITH) considers continued growth directed inward. Virtual worlds may contain agents or processes that generate further worlds, producing an approximately self-similar hierarchy of designed environments, institutions, histories, and experiences. A civilization need not become passive or lose its ambition. It may increase the quantity, variety, or quality of valued activity without repeatedly establishing resident infrastructure elsewhere. Technological progress may also make the supporting services of a chosen internal growth path sufficiently available within existing territory. The strong hypothesis combines these possibilities: \emph{continued internal growth can coexist with no additional external investment when local support remains sufficient and the relevant allocation inequalities continue to hold.} Recursively nested worlds motivate a workload-accounting example; the allocation mechanism does not require recursion.

This is not the assertion that an ordinary finite computer implements infinitely many independent worlds at once. Nor is a computer inside a simulation an additional physical computer that can be added to the host's capacity. \CITH\ is a production and allocation hypothesis subject to physical accounting, not a source of free substrate. Energy, state storage, entropy management, replacement material, error correction, and failure risk remain relevant \citep{Landauer1961,Bennett1973,Lloyd2000}. The distinction between a generative rule and its instantiated states will be explicit.

\subsection{Relation to earlier work and the claimed contribution}

The qualitative ideas of postbiological civilizations and inward development are not new. \citet{Dick2003} and \citet{CirkovicBradbury2006} considered computationally oriented civilizations; \citet{Smart2012} proposed transcension toward increasingly dense internal scales. Aestivation concerns the timing of computation rather than substituting internal development for external investment \citep{Sandberg2017}. The thermodynamic critique by \citet{BennettHansonRiedel2019} is important here: efficiency cannot be evaluated without specifying accessible entropy reservoirs. Sustainability and homeostasis offer different mechanisms for bounded growth \citep{HaqqMisraBaum2009,WongBartlett2022}. The present paper neither imports a universal evolutionary destiny nor infers preferences from energy efficiency alone.

A non-peer-reviewed working paper uses the phrase ``Inward Transition'' for a post-AGI shift toward internal computation and reduced detectability \citep{Pitti2026}. No priority is claimed for that phrase, virtual worlds, or inwardness. Conversely, resource acquisition and self-replicating industry can strengthen expansion incentives, including on intergalactic scales \citep{ArmstrongSandberg2013}; autonomous machine industry may reduce dependence on biological settlement \citep{Ivliev2026}. Population-total value can make unused astronomical resources represent substantial foregone welfare \citep{Bostrom2003Waste}. These are direct counterpressures, not exceptions to be excluded by definition. Social-scientific critiques also caution against treating poorly identified motives or human scaling laws as universal civilizational laws \citep{Shkurko2024,JacksonCriadoPerez2024}.

The proposed contribution is a more specific combination: a workload frontier, a common-budget comparison of internal and external activity, exact logarithmic thresholds, a curvature trichotomy, and pathwise persistence conditions. Access costs, larger budgets, residual external value, production curvature, and future resource production supply explicit tests of the mechanism. The optimization and branching-process tools are standard. Their application identifies assumptions that a strong inward-growth claim must satisfy and circumstances in which it fails. The analytical thresholds, rather than a fraction of a synthetic ensemble, are the main result.

\subsection{Four distinct claims}

The strong inward-growth claim must be separated from three neighbouring statements. First, fewer new resident settlements may coexist with extensive external mining and probes; this is a weaker \emph{settlement-composition} branch. Second, either branch could be bright or faint. Third, even common inward allocation need not explain galactic silence because autonomous expansionist descendants can dominate territory. We distinguish the \emph{allocation claim}, the \emph{detectability claim}, and the \emph{population claim}. The formal focus is all-inward allocation, not merely reduced settlement. Neither detectability nor a population-level explanation of silence follows from that allocation alone.

